% =====================================================================
%  Submission to AI4Nature Workshop @ AVSS 2026 (Lecce, Italy)
%  Format: IEEE Conference Proceedings (double column, up to 6 pages)
%
%  Compile with pdfLaTeX on Overleaf.
%  Required files in project root:
%     main.tex, ref.bib
%     figs/fig_world.png
%     figs/fig_reliability.png
%     figs/fig_per_cell_estimates.png
%     figs/fig_allocations.png
%     figs/fig_regret_vs_T.png
%     figs/fig_regret_vs_calsize.png
%     figs/fig_coverage.png
% =====================================================================

\documentclass[conference]{IEEEtran}
\IEEEoverridecommandlockouts

\usepackage{cite}
\usepackage{amsmath,amssymb,amsfonts}
\usepackage{algorithm,algpseudocode}
\usepackage{graphicx}
\usepackage{textcomp}
\usepackage{xcolor}
\usepackage{booktabs}
\usepackage{siunitx}
\usepackage{url}
\usepackage{hyperref}
\hypersetup{colorlinks=true,linkcolor=black,citecolor=black,urlcolor=blue}

\def\BibTeX{{\rm B\kern-.05em{\sc i\kern-.025em b}\kern-.08em
    T\kern-.1667em\lower.7ex\hbox{E}\kern-.125emX}}

\begin{document}

% ----- Title ---------------------------------------------------------
\title{From Pixels to Patrols:\\
A Calibration-Aware Sensor Fusion Pipeline\\
for Camera-Trap-Driven Anti-Poaching\\
Resource Allocation}

% ----- Authors -------------------------------------------------------
%\author{
%\IEEEauthorblockN{Ka Hin Chan}
%\IEEEauthorblockA{Col\'egio Diocesano \\
%de S\~ao Jos\'e (5\textordfeminine)\\
%Macao SAR, China}
%\and
%\IEEEauthorblockN{Long Nam Ao}
%\IEEEauthorblockA{Col\'egio Diocesano \\
%de S\~ao Jos\'e (5\textordfeminine)\\
%Macao SAR, China}
%\and
%\IEEEauthorblockN{Weng Kin Loi}
%\IEEEauthorblockA{Col\'egio Diocesano \\
%de S\~ao Jos\'e (5\textordfeminine)\\
%Macao SAR, China}
%\and
%\IEEEauthorblockN{Long Tin Tse}
%\IEEEauthorblockA{Col\'egio Diocesano \\
%de S\~ao Jos\'e (5\textordfeminine)\\
%Macao SAR, China}
%\and
%\IEEEauthorblockN{Sok Kin Cheng\textsuperscript{\dag}}
%\IEEEauthorblockA{Col\'egio Diocesano \\
%de S\~ao Jos\'e (5\textordfeminine)\\
%Macao SAR, China \\
%\texttt{chengsokkin@cdsj5.edu.com}}
%}

\maketitle

\renewcommand{\thefootnote}{\dag}
\footnotetext{Corresponding author.}
\renewcommand{\thefootnote}{\arabic{footnote}}

% ----- Abstract ------------------------------------------------------
\begin{abstract}
Camera-trap networks coupled with deep image classifiers now generate
species-level detections at continental scale, but operational
anti-poaching decisions taken downstream of these classifiers typically
treat the predictions as ground truth. We argue that this discards the
detection uncertainty that ecologists have begun to quantify, and that
the discard is operationally costly. Building on the recent
demonstration that camera-trap classifiers are systematically
miscalibrated (Dussert et al., 2025), we propose an end-to-end
calibration-aware pipeline that propagates classifier uncertainty
through to spatial patrol allocation. The pipeline composes (i)~post-hoc
temperature scaling, (ii)~split-conformal regression on the per-cell
expected threat count, and (iii)~four allocation policies
(\emph{oracle}, \emph{naive}, \emph{calibrated point estimate},
\emph{distributionally robust upper}) on a water-filling objective. On a
controlled synthetic camera-trap network of $400$ cells and roughly
$16{,}000$ images calibrated to the over-confidence regime reported in
the literature, temperature scaling reduces expected calibration error
from $0.34$ to $0.05$ and closes $79\%$ of the decision-stage regret of
the naive policy; the conformal-set distributionally robust policy
achieves comparable regret while maintaining empirical coverage near
the nominal $90\%$ level. The cost of decision-stage miscalibration grows steeply
in the over-confident regime and more mildly in the under-confident
regime, and the conformal interval's coverage tracks its nominal value
as the calibration sample size grows. The framework slots directly under
existing camera-trap classifiers (MegaDetector, SpeciesNet) and is
substitutable for the kernel-density occurrence priors used in current
spatial planning systems.
\end{abstract}

\begin{IEEEkeywords}
Camera traps, calibration, conformal prediction, anti-poaching, patrol
allocation, sensor fusion, distributionally robust optimization.
\end{IEEEkeywords}

% =====================================================================
\section{Introduction}
% =====================================================================

The deep-learning revolution has transformed camera-trap ecology.
MegaDetector~\cite{morris2025megadetector,gadot2024crop},
SpeciesNet~\cite{gadot2024crop}, and an ecosystem of regional
classifiers~\cite{norouzzadeh2018automatically} now process hundreds of
millions of images per year, producing per-image species predictions
that feed downstream conservation tasks: occupancy estimation, species
interaction modelling, real-time poaching alerts~\cite{whytock2023real},
and---of central interest here---spatial planning of ranger patrols.

The downstream pipeline almost universally treats classifier outputs as
ground truth. Yet a recent and operationally significant finding by
Dussert et al.~\cite{dussert2025confident} is that camera-trap
classifiers are systematically \emph{miscalibrated}: their softmax
confidences do not correspond to empirical accuracies. The authors
report large expected calibration errors and argue that the bias in
confidence scores then biases the inferential tasks downstream of the
classifier. Their proposed remedy is post-hoc
calibration~\cite{guo2017calibration,platt1999probabilistic}, and their
empirical investigation is restricted to \emph{inferential} downstream
tasks (occupancy estimation, interaction inference, real-time
alerting). The implicit conjecture is that the calibration question is
inferential in nature.

We argue otherwise: calibration is also a \emph{decision-stage}
question, and one that has been overlooked. When per-cell species
detections feed into a spatial resource allocation, miscalibration
distorts the relative ranking of cells, which distorts the allocation,
which produces measurable regret against an oracle that knows the true
threat field. The cost is not merely cosmetic: in anti-poaching
operations, mis-allocated effort has direct welfare consequences for
endangered species. To our knowledge, no existing work in either the
calibration literature or the security-game literature
(PAWS~\cite{fang2016paws}, CAPTURE~\cite{nguyen2016capture},
distributionally robust
optimisation~\cite{esfahani2018wasserstein}) propagates classifier
calibration explicitly through to spatial allocation.

\noindent\textbf{Contributions.} We propose an end-to-end pipeline that
makes this propagation explicit and quantifies the resulting gains:
\begin{itemize}
  \item We formulate camera-trap-driven patrol allocation as a
    \emph{sensor fusion} problem: per-image softmax outputs are
    temperature-scaled, aggregated to per-cell expected counts of the
    threat class, wrapped in a split-conformal regression interval, and
    fed into a budget-constrained water-filling allocator.
  \item We define four allocation policies on this pipeline---oracle,
    naive (uncalibrated point estimate), calibrated (temperature-scaled
    point estimate), and distributionally robust (conformal
    upper)---and an oracle-relative regret metric.
  \item In a controlled synthetic experiment calibrated to the
    over-confidence regime reported by~\cite{dussert2025confident}, we
    show that calibration reduces decision regret from $3.0\%$ to
    $0.63\%$ of the oracle utility---closing $79\%$ of the naive policy
    gap---and that conformal intervals achieve nominal coverage in the
    large-calibration-set regime.
\end{itemize}

% =====================================================================
\section{Related Work}
% =====================================================================

\textbf{Camera-trap computer vision.} Norouzzadeh et
al.~\cite{norouzzadeh2018automatically} initiated the modern era of
deep-learning species classification on Snapshot Serengeti; recent
benchmarks combine MegaDetector~\cite{morris2025megadetector} with
classifiers such as SpeciesNet~\cite{gadot2024crop}, and out-of-domain
generalisation remains a recognised challenge~\cite{norman2023cnn,
beery2018recognition,pahuja2024cosmo}. Camera-trap-driven real-time
alerts have been deployed in
Gabon~\cite{whytock2023real}.

\textbf{Calibration in conservation.} Outside the conservation
literature, post-hoc calibration of neural networks is by now standard:
temperature scaling~\cite{guo2017calibration} and its
extensions~\cite{platt1999probabilistic}. Dussert et
al.~\cite{dussert2025confident} are the first systematic treatment of
the question for camera-trap classifiers, reporting expected
calibration errors of comparable magnitude to those produced by our
synthetic generator. Their downstream evaluation, however, is on
inferential tasks; the present paper extends the analysis to
decision-stage tasks.

\textbf{Conformal prediction.} The split-conformal framework of
Vovk et al.~\cite{vovk2005algorithmic}---reviewed accessibly
in~\cite{angelopoulos2023conformal}---provides distribution-free
finite-sample coverage guarantees on prediction sets. The most common
applications are to classification~\cite{romano2020classification};
applications to ecological prediction are scarce, and to our knowledge
none integrate with downstream spatial allocation.

\textbf{Patrol allocation and Stackelberg models.} Spatial planning of
ranger patrols originated with PAWS~\cite{fang2016paws} and
CAPTURE~\cite{nguyen2016capture}; a controlled field
trial~\cite{gholami2017taking} demonstrated operational benefit.
Distributionally robust optimisation~\cite{esfahani2018wasserstein}
provides the mathematical machinery to hedge against ambiguity in the
input distribution, and we use a conformal-set instantiation of this
idea below.

% =====================================================================
\section{Method}
% =====================================================================

\subsection{Problem and Notation}

The protected area is partitioned into $N$ cells. At each cell $i\in
\{1,\dots,N\}$ a camera trap collects $M_i$ images; the classifier
emits per-image probability vectors
$\boldsymbol{\pi}_{i,j}\in\Delta^{K-1}$ over $K$ species classes,
where class~$K{-}1$ is the threat class of interest (e.g.\ the
human/poacher signature class returned by MegaDetector). The
defender allocates a total budget $T$ of patrol-hours over the cells,
choosing $\tau_i\ge0$ with $\sum_i\tau_i=T$ and $\tau_i\le\tau_{\max}$.
The capture utility from cell~$i$ is
\begin{equation}
  u_i(\tau_i,z_i) \;=\; z_i\bigl(1-e^{-\lambda\,\tau_i}\bigr),
  \label{eq:cell_utility}
\end{equation}
where $z_i$ is the true expected count of threat events in cell~$i$
and $\lambda$ a known detection efficiency. The aggregate utility is
$U(\boldsymbol{\tau};\mathbf{z})=\sum_i u_i(\tau_i,z_i)$, which is
concave and separable in $\boldsymbol{\tau}$. The decision-maker
observes only $\{\boldsymbol{\pi}_{i,j}\}$, not $\mathbf{z}$.

\subsection{Calibration}

A small fraction $\rho_{\mathrm{cal}}$ of the images is held out with
labels, yielding a calibration set
$\mathcal{C}=\{(\boldsymbol{\pi}_j,y_j)\}_{j=1}^{n_{\mathrm{cal}}}$.
Following~\cite{guo2017calibration,dussert2025confident} we recover
logits up to translation, $\ell_j=\log\boldsymbol{\pi}_j$, and fit
\begin{equation}
  \hat{T}\;=\;\arg\min_{T>0}\;\;-\!\sum_{j\in\mathcal{C}}
  \log\Bigl[\,\mathrm{softmax}(\ell_j/T)\,\Bigr]_{y_j},
  \label{eq:temperature}
\end{equation}
and set $\tilde{\boldsymbol{\pi}}_{i,j}=\mathrm{softmax}(\ell_{i,j}/\hat T)$
on deployment. Per-cell calibrated expected counts of the threat class
are
\begin{equation}
  \hat{z}_i \;=\; \sum_{j\in\mathcal{D}_i}\tilde{\boldsymbol{\pi}}_{i,j,K-1},
  \label{eq:zhat}
\end{equation}
where $\mathcal{D}_i$ indexes deployment images at cell~$i$.

\subsection{Split-Conformal Regression on Per-Cell Counts}

Standardised residuals are computed cell-wise on the calibration
images. For cell $i$ with calibration count
$n^{\mathrm{cal}}_i\ge n_{\min}$, write
$z^{\mathrm{tr}}_i=\sum_{j\in\mathcal{C}_i}\mathbf{1}\{y_j=K{-}1\}$
and $\hat{z}^{\mathrm{cal}}_i=\sum_{j\in\mathcal{C}_i}\tilde{\boldsymbol{\pi}}_{i,j,K-1}$.
Define
\begin{equation}
  r_i \;=\; (z^{\mathrm{tr}}_i - \hat{z}^{\mathrm{cal}}_i)/\sqrt{n^{\mathrm{cal}}_i}.
\end{equation}
The $\sqrt{n^{\mathrm{cal}}_i}$ standardisation is motivated by
$\mathrm{Var}\bigl(\sum_{j}\mathbf{1}\{y_j=K{-}1\}\bigr)\propto n^{\mathrm{cal}}_i$
for cell-level Bernoulli sums and makes the residuals approximately
exchangeable across heterogeneous cell sizes. Taking the $(1-\alpha)$
empirical quantile $q$ of $\{|r_i|\}$, the per-cell prediction interval
on deployment is
\begin{equation}
  I_i \;=\; \bigl[\,\hat{z}_i - q\sqrt{n^{\mathrm{dep}}_i},\;
                       \hat{z}_i + q\sqrt{n^{\mathrm{dep}}_i}\bigr]\cap[0,n^{\mathrm{dep}}_i].
  \label{eq:interval}
\end{equation}
By the standard split-conformal argument and exchangeability of cells,
$\mathbb{P}(z_i^{\star}\in I_i)\ge 1-\alpha$ marginally.

\subsection{Four Allocation Policies}

The allocator solves
$\max_{\boldsymbol{\tau}}\;\sum_i z_i^{\mathrm{plan}}(1-e^{-\lambda\tau_i})$
subject to $\sum_i\tau_i=T$ and $0\le\tau_i\le\tau_{\max}$, which by
Karush--Kuhn--Tucker gives the water-filling closed form
$\tau_i^{\star}(\nu)=\mathrm{clip}_{[0,\tau_{\max}]}\!\bigl(\lambda^{-1}\log(\lambda z_i^{\mathrm{plan}}/\nu)\bigr)$
with $\nu$ the dual variable; we solve for $\nu$ by bisection. The four
policies differ only in the planning value $z_i^{\mathrm{plan}}$:
\begin{itemize}
  \item \emph{Oracle:} $z_i^{\mathrm{plan}} = z_i^{\star}$,
        evaluated for benchmarking.
  \item \emph{Naive:} $z_i^{\mathrm{plan}} = \sum_j\boldsymbol{\pi}_{i,j,K-1}$,
        uncalibrated soft count.
  \item \emph{Calibrated:} $z_i^{\mathrm{plan}}=\hat{z}_i$ from
        \eqref{eq:zhat}.
  \item \emph{DRO (conformal upper):}
        $z_i^{\mathrm{plan}}=\hat{z}_i+q\sqrt{n^{\mathrm{dep}}_i}$, the
        upper endpoint of $I_i$. This realises a risk-averse policy
        that hedges against the false-negative case in which true
        threat is larger than the calibrated point estimate---the
        relevant asymmetry for security
        applications~\cite{esfahani2018wasserstein}.
\end{itemize}
For policy~$P$ the decision-stage \emph{regret} is
$R_P=U(\boldsymbol{\tau}^{\star};\mathbf{z}^{\star})
-U(\boldsymbol{\tau}^{P};\mathbf{z}^{\star})$,
evaluated under the true threat field.

% =====================================================================
\section{Numerical Experiments}
% =====================================================================

\subsection{Synthetic Camera-Trap Network}

The model is evaluated on a $20\times 20$ grid ($N=400$) with $K=4$
classes; the threat class has spatial structure drawn from a Fourier-filtered
Gaussian random field with correlation length $\rho=3$ cells, and
$M_i\sim\mathrm{Poisson}(40)$ images per cell, totalling roughly
$16{,}000$ images.  The classifier injects logits
$\ell=\log p_i+\mathbf{b}+\sigma_{\mathrm{obs}}\boldsymbol{\varepsilon}$
where $\mathbf{b}$ is a per-class bias vector with positive entry on
the threat class, $\sigma_{\mathrm{obs}}=0.8$, and emits
$\boldsymbol{\pi}=\mathrm{softmax}(\ell/T_{\mathrm{true}})$ with
$T_{\mathrm{true}}=0.45$.  This is the \emph{over-confident} regime
reported by~\cite{dussert2025confident} for deep camera-trap
classifiers: the bias inflates the threat-class log-probability while
the small temperature sharpens distributions toward the inflated mode.
We allocate $T=60$ patrol units, $\tau_{\max}=3$, $\lambda=0.25$, and
use $\rho_{\mathrm{cal}}=0.20$, $\alpha=0.10$, $n_{\min}=3$. The
ground-truth threat field, image density, and true expected threat
count appear in Fig.~\ref{fig:world}.

\begin{figure*}[t]
  \centering
  \includegraphics[width=0.96\linewidth]{figs/fig_world.png}
  \caption{The synthetic camera-trap network. Left: true
  $P(\mathrm{threat})$ per cell, spatially correlated via
  Fourier-filtered Gaussian random field. Middle: realised number of
  images per cell, Poisson with mean 40. Right: true expected threat
  count $z_i^{\star}$ per cell, the quantity allocation must
  estimate.}
  \label{fig:world}
\end{figure*}

\subsection{Calibration}

Fig.~\ref{fig:rel} shows the reliability diagrams before and after
temperature scaling on the calibration set. The uncalibrated classifier
is strongly over-confident with expected calibration error
$\mathrm{ECE}=0.344$ and multi-class Brier score $0.852$; after fitting
$\hat T=4.27$ via~\eqref{eq:temperature}, $\mathrm{ECE}$ falls to
$0.052$ and Brier to $0.660$.  Both the magnitude and the direction of
the calibration improvement are consistent with the empirical results
reported by~\cite{dussert2025confident} on real
camera-trap data, which lends our synthetic generator face validity for
the decision-stage analysis that follows.

\begin{figure}[t]
  \centering
  \includegraphics[width=0.98\linewidth]{figs/fig_reliability.png}
  \caption{Reliability diagrams before (left) and after (right)
  temperature scaling. The uncalibrated classifier is over-confident
  (mean confidence systematically above accuracy); the scaled
  classifier closely tracks the diagonal.}
  \label{fig:rel}
\end{figure}

\subsection{Per-Cell Estimates and Conformal Coverage}

Fig.~\ref{fig:percell} compares naive and calibrated per-cell threat
estimates against the deployment-subset truth, and shows the empirical
coverage of the conformal intervals from~\eqref{eq:interval}. The naive
estimator over-shoots truth (linear-fit slope of $1.40$ relative to the
$45^{\circ}$ line); the calibrated estimator tracks truth closely (slope
$0.97$). Across $400$ cells the conformal interval contains the true
$z_i^{\star}$ in $84.5\%$ of cells at the nominal $90\%$ level, slightly
under-covering at this calibration size; the coverage--size sweep below
shows the asymptotic guarantee is recovered as $\rho_{\mathrm{cal}}$
grows.

\begin{figure*}[t]
  \centering
  \includegraphics[width=0.96\linewidth]{figs/fig_per_cell_estimates.png}
  \caption{Per-cell estimates and conformal coverage. Left: naive
  (uncalibrated) estimates overshoot. Middle: calibrated estimates
  track truth. Right: the $90\%$ conformal band (green) brackets the
  true per-cell expected threat count (black) over $400$ cells sorted
  by truth.}
  \label{fig:percell}
\end{figure*}

\subsection{Allocations}

Fig.~\ref{fig:alloc} maps the four allocations. The naive policy
allocates patrols to inflated false hotspots, leaving genuine hotspots
underprotected; the calibrated and DRO allocations track the oracle
closely. Quantitatively, on the true threat field
\begin{itemize}
  \item Oracle utility $U^{\star}=373.6$.
  \item Naive utility $362.4$; regret $11.19$ ($2.99\%$ of $U^{\star}$).
  \item Calibrated utility $371.3$; regret $2.35$ ($0.63\%$).
  \item DRO utility $370.7$; regret $2.88$ ($0.77\%$).
\end{itemize}
Calibration alone closes $\mathbf{78.96\%}$ of the naive policy's
regret gap; DRO closes $\mathbf{74.25\%}$ at the same nominal
calibration size.  The remaining $\sim$$0.6\%$ gap is irreducible
estimation variance from the finite calibration set.

\begin{figure*}[t]
  \centering
  \includegraphics[width=0.96\linewidth]{figs/fig_allocations.png}
  \caption{Allocation maps on the $20\times 20$ grid. The oracle
  concentrates rangers on the true threat hotspots; the naive
  allocation diverts rangers to false hotspots induced by classifier
  over-confidence; the calibrated and DRO allocations recover the
  oracle's hotspot structure.}
  \label{fig:alloc}
\end{figure*}

\subsection{Sensitivity to Miscalibration Severity}

Fig.~\ref{fig:Tsweep} sweeps $T_{\mathrm{true}}$ from heavy
over-confidence ($T_{\mathrm{true}}=0.30$) through well-calibrated
($T_{\mathrm{true}}=1.0$) to heavy under-confidence
($T_{\mathrm{true}}=2.50$), averaging four random seeds per point. The
naive policy's regret grows steeply in the over-confident regime and
mildly in the under-confident regime, with a local minimum near
$T_{\mathrm{true}}\!\approx\!1.2$ where temperature softening cancels
the additive class bias.  The calibrated and DRO policies are
direction-agnostic and stay at $\sim$$0.6$--$0.7\%$ regret across the
entire sweep, dominating the naive policy at every $T_{\mathrm{true}}\le 1$
tested---which is the operationally relevant regime since modern deep
classifiers are typically over-confident~\cite{guo2017calibration,dussert2025confident}.

\begin{figure}[t]
  \centering
  \includegraphics[width=0.98\linewidth]{figs/fig_regret_vs_T.png}
  \caption{Decision-stage regret versus miscalibration severity. The
  naive policy's regret grows with miscalibration magnitude;
  calibrated and DRO policies are flat. The current empirical regime
  for camera-trap classifiers~\cite{dussert2025confident} sits in
  the over-confident half ($T_{\mathrm{true}}<1$).}
  \label{fig:Tsweep}
\end{figure}

\subsection{Calibration Sample Size and Conformal Coverage}

Fig.~\ref{fig:calsweep} reports regret and coverage as
$\rho_{\mathrm{cal}}$ varies from $0.05$ to $0.60$, averaged over five
seeds. The calibrated policy's regret is small and approximately flat
in $\rho_{\mathrm{cal}}$, the DRO policy's regret grows with
$\rho_{\mathrm{cal}}$ because larger calibration sets shrink the
deployment set and reduce per-cell sample size for the soft-count
estimator, and empirical coverage tracks the nominal $90\%$ level
closely from $\rho_{\mathrm{cal}}=0.20$ onwards
(Fig.~\ref{fig:coverage}). The under-coverage at small
$\rho_{\mathrm{cal}}$ is a finite-sample artifact of the cell-level
exchangeability assumption and is consistent with the standard
split-conformal asymptotic
analysis~\cite{angelopoulos2023conformal}.

\begin{figure}[t]
  \centering
  \includegraphics[width=0.98\linewidth]{figs/fig_regret_vs_calsize.png}
  \caption{Regret versus calibration set size. The naive policy
  regret is approximately constant (calibration set is unused); the
  calibrated policy regret is small throughout; the DRO regret grows
  modestly because deployment-set size shrinks correspondingly.}
  \label{fig:calsweep}
\end{figure}

\begin{figure}[t]
  \centering
  \includegraphics[width=0.84\linewidth]{figs/fig_coverage.png}
  \caption{Empirical coverage of the conformal interval versus
  calibration set fraction. Coverage approaches the nominal $90\%$
  from below as $\rho_{\mathrm{cal}}$ grows, as predicted by the
  finite-sample split-conformal bound.}
  \label{fig:coverage}
\end{figure}

% =====================================================================
\section{Discussion}
% =====================================================================

Three points follow from the numerical experiments.

\emph{First}, decision-stage miscalibration is operationally costly.
For a classifier with the over-confidence profile reported on real
camera-trap data by~\cite{dussert2025confident}, the naive policy
loses approximately $3\%$ of oracle utility---a non-trivial loss when
total interception capacity is binding. The cost is steepest in the
over-confident regime---the regime reported for deployed DL
classifiers---and tapers off as miscalibration shrinks; under-confident
classifiers are penalised similarly in direction because the relative
ranking of cells is the operational quantity, not the absolute scale.

\emph{Second}, the remedy is cheap. Temperature scaling has one scalar
parameter and is fit in $O(n_{\mathrm{cal}})$ time on an existing
calibration split that any deployment of a camera-trap classifier
already produces for accuracy auditing. The conformal layer adds a
single quantile computation and an empirically validated coverage
guarantee. The pipeline therefore composes additively under existing
camera-trap classifiers and adds negligibly to inference time.

\emph{Third}, distributionally robust optimisation against the
conformal upper bound buys robustness when the operational asymmetry
favours false alarms over missed threats---the standard regime for
anti-poaching. In our experiment the DRO policy approximately matches
the calibrated point policy on regret, while planning against intervals
rather than points; in deployments with operationally tractable
asymmetry costs, the DRO formulation can be re-weighted to trade off
regret against false-negative rate.

\textbf{Limitations.} The principal limitation is the synthetic nature
of the experiment. Replacing the synthetic classifier outputs with
real MegaDetector/SpeciesNet inferences on a held-out camera-trap
benchmark such as iWildCam~\cite{beery2020iwildcam} or Snapshot
Serengeti is the most direct extension, and we have designed the
pipeline so that the substitution is module-local: the classifier in
\S{}IV-A becomes the real model, no other component changes. A second
limitation is the marginal-coverage guarantee of split conformal,
which does not provide conditional coverage on, e.g., subpopulations
of cells with rare habitats. Cluster-conditional or class-conditional
conformal extensions are available, at the cost of additional
calibration data. Third, the synthetic threat field is spatially
correlated but temporally static; a real-time pipeline would need to
re-fit $\hat T$ periodically as the deployment regime drifts.

% =====================================================================
\section{Conclusion}
% =====================================================================

Operational use of camera-trap classifiers for anti-poaching patrol
allocation is, today, an inferential pipeline read as if it were a
decision pipeline. We have proposed and evaluated a small,
self-contained set of additions---temperature scaling, split-conformal
regression on per-cell expected threat count, and four candidate
allocation policies---that close approximately four-fifths of the
decision-stage regret incurred by naive uncalibrated allocation, while
preserving the existing classifier untouched. The components are
substitutable into existing camera-trap pipelines and into existing
patrol-planning systems (PAWS, CAPTURE). The natural next step is to
apply the framework on real classifier outputs from a deployed
camera-trap network, which would convert it from a methodological
demonstration into a calibrated operational tool. The accompanying
implementation is open source at
\url{https://github.com/cdsj5-ai4nature/calibration-aware-patrol}.

% =====================================================================
\section*{Acknowledgment}
% =====================================================================

The authors thank Col\'egio Diocesano de S\~ao Jos\'e
(5\textordfeminine) in Macao SAR for supporting this work.

\bibliographystyle{IEEEtran}
\bibliography{ref}

\end{document}
